\documentclass[11pt]{article}

\usepackage{amsmath}
\usepackage{amsfonts}
\usepackage{graphicx}
\usepackage{hyperref}

\title{MedVision AI \\ Mathematical and Algorithmic Foundations}
\author{Yann Smatti}
\date{}

\begin{document}

\maketitle

\tableofcontents

\newpage

\section{Introduction}

Medical image analysis has become a major application of machine learning and computer vision. 
Modern diagnostic systems assist radiologists by detecting patterns in medical images such as:

\begin{itemize}
\item X-Ray scans
\item CT scans
\item MRI images
\end{itemize}

This document explains the mathematical foundations of the algorithms implemented in the MedVision AI project.

\section{Image Representation}

An image is mathematically represented as a matrix.

For a grayscale image:

\[
I \in \mathbb{R}^{H \times W}
\]

where

\begin{itemize}
\item $H$ = height
\item $W$ = width
\end{itemize}

Each pixel value represents intensity:

\[
I(x,y) \in [0,255]
\]

After normalization:

\[
I(x,y) \in [0,1]
\]

\section{Convolutional Neural Networks}

CNNs are the standard architecture for image analysis.

A convolution operation is defined as:

\[
(I * K)(x,y) = \sum_i \sum_j I(x-i,y-j)K(i,j)
\]

where

\begin{itemize}
\item $I$ is the image
\item $K$ is the convolution kernel
\end{itemize}

Example kernel:

\[
K =
\begin{bmatrix}
1 & 0 & -1 \\
1 & 0 & -1 \\
1 & 0 & -1
\end{bmatrix}
\]

Convolution allows feature extraction such as:

\begin{itemize}
\item edges
\item textures
\item shapes
\end{itemize}

\section{Activation Functions}

Neural networks use nonlinear activation functions.

The most common is ReLU:

\[
ReLU(x) = \max(0,x)
\]

This introduces non-linearity and improves training stability.

\section{Pooling}

Pooling reduces spatial dimensions.

Max pooling:

\[
y = \max(x_1,x_2,...,x_n)
\]

Benefits:

\begin{itemize}
\item reduces computation
\item improves invariance
\item reduces overfitting
\end{itemize}

\section{Loss Functions}

Binary classification uses Binary Cross Entropy:

\[
L = -\frac{1}{N} \sum_{i=1}^{N} [y_i \log(p_i) + (1-y_i)\log(1-p_i)]
\]

where

\begin{itemize}
\item $y_i$ = true label
\item $p_i$ = predicted probability
\end{itemize}

\section{Transfer Learning}

Training deep networks from scratch requires large datasets.

Transfer learning uses pretrained models such as:

\begin{itemize}
\item ResNet
\item EfficientNet
\item VGG
\end{itemize}

Mathematically:

\[
f(x) = g(h(x))
\]

where

\begin{itemize}
\item $h(x)$ = pretrained feature extractor
\item $g(x)$ = task specific classifier
\end{itemize}

\section{Evaluation Metrics}

Medical models require robust evaluation metrics.

Accuracy:

\[
Accuracy = \frac{TP + TN}{TP + TN + FP + FN}
\]

Precision:

\[
Precision = \frac{TP}{TP + FP}
\]

Recall (Sensitivity):

\[
Recall = \frac{TP}{TP + FN}
\]

F1 Score:

\[
F1 = 2 \cdot \frac{Precision \cdot Recall}{Precision + Recall}
\]

\section{ROC Curve}

The ROC curve represents the tradeoff between

\begin{itemize}
\item True Positive Rate
\item False Positive Rate
\end{itemize}

\[
TPR = \frac{TP}{TP + FN}
\]

\[
FPR = \frac{FP}{FP + TN}
\]

\section{Explainable AI}

In medical applications, explainability is crucial.

Grad-CAM computes:

\[
L^c = ReLU\left(\sum_k \alpha_k^c A^k\right)
\]

where

\begin{itemize}
\item $A^k$ are feature maps
\item $\alpha_k^c$ are importance weights
\end{itemize}

This produces a heatmap highlighting regions influencing the prediction.

\section{Conclusion}

Computer vision combined with deep learning provides powerful tools for medical diagnosis assistance.

However, rigorous validation and explainability remain essential.

\end{document}